\documentclass[11pt,letterpaper]{ctexart}

\usepackage[margin=1.35in,top=1.15in,bottom=1.0in,headheight=14pt]{geometry}
\usepackage{amsmath,amssymb}
\usepackage[ruled,vlined,linesnumbered]{algorithm2e}
\usepackage{fancyhdr}
\usepackage{hyperref}

\hypersetup{
  colorlinks=false,
  pdfauthor={你的名字},
  pdftitle={基于强化学习与蒙特卡洛树搜索的量子布尔函数综合方法}
}

\pagestyle{fancy}
\fancyhf{}
\fancyhead[L]{\nouppercase{\leftmark}}
\fancyhead[R]{\thepage}
\renewcommand{\headrulewidth}{0pt}

\fancypagestyle{plain}{
  \fancyhf{}
  \fancyfoot[C]{\thepage}
  \renewcommand{\headrulewidth}{0pt}
}

\SetKwInput{KwInput}{Input}
\SetKwInput{KwOutput}{Output}
\SetKwComment{tcp}{// }{}
\DontPrintSemicolon
\SetAlgoNlRelativeSize{0}
\SetAlgoInsideSkip{smallskip}

\begin{document}

\thispagestyle{plain}
\begin{center}
\vspace*{1.35in}
{\LARGE 基于强化学习与蒙特卡洛树搜索的量子布尔\\[0.35em]函数综合方法\par}
\vspace{0.42in}
{\large 你的名字\par}
\vspace{0.22in}
{\large 2026 年 7 月 1 日\par}
\end{center}
\vspace{0.42in}

\section{方法}

本文提出的方法将量子布尔函数综合建模为一个 ** 序列决策问题 **，
并采用强化学习（RL）与蒙特卡洛树搜索（MCTS）相结合的智能框架进行
求解。

\subsection{问题形式化}

给定一个 $n$ 变量布尔函数的真值表 $T \in \{0,1\}^{2^n}$，目标为找到一个等价
的量子电路，即一个门序列 $G=(g_1,g_2,\ldots,g_L)$，使得对应的酉矩阵 $U(G)$
满足 $U(G)|0\rangle^{\otimes n}=|T\rangle$（其中 $|T\rangle$ 为真值表对应的计算基态叠加）。

我们将该问题建模为 ** 马尔可夫决策过程（MDP）**，定义如下：

\begin{itemize}
  \item \textbf{状态 $s_t$：} 当前部分构建的电路状态，包括目标真值表 $T$ 和当前电路的
  酉矩阵 $U_t$。

  \item \textbf{动作 $a_t$：} 从预定义动作空间 $\mathcal{A}$ 中选择一个量子门（含门类型和作用比
  特），即 $a_t=(gate,qubits)$。

  \item \textbf{状态转移 $s_{t+1}=s_t\circ a_t$：} 即在当前电路后添加该门。

  \item \textbf{奖励 $R(s_t,a_t)$：} 当电路构建完成（达到终止状态）时，根据最终电路的
  保真度 $\mathcal{F}$ 给予稀疏奖励：
  \[
    R_{\mathrm{final}}=
    \begin{cases}
      +1, & \mathcal{F}=1 \quad \text{(电路完全正确)},\\
      -1, & \mathcal{F}<1 \quad \text{(电路错误)},
    \end{cases}
  \]
  并可附加过程惩罚 $r_{\mathrm{step}}=-\lambda$ 以鼓励简洁电路。

  \item \textbf{终止条件：} 电路深度达到预设最大值，或电路已满足目标函数。
\end{itemize}

\subsection{双网络智能体架构}

智能体由策略网络和价值网络构成，二者共享一个主干编码器，但输出
头不同。

\subsubsection{状态表示}

输入状态由两部分拼接而成：

\begin{enumerate}
  \item 目标真值表 $T$ 的向量化表示（长度 $2^n$）；

  \item 当前酉矩阵 $U_t$ 的实部与虚部展平（维度 $2 \times 2^n \times 2^n$）。
\end{enumerate}

为避免维度过大，可在实际实现中对酉矩阵进行降维（如 PCA 或自编码
器），本文采用直接展平方式以保持信息完整。

\subsubsection{策略网络 $\pi_\theta(a|s)$}

策略网络输出一个概率分布，对每个动作 $a \in \mathcal{A}$ 打分：
\[
  \pi_\theta(a|s)=\operatorname{Softmax}(f_{\mathrm{policy}}(h_s)),
\]
其中 $h_s$ 为共享主干提取的特征向量，$f_{\mathrm{policy}}$ 为若干全连接层。该输出作为
MCTS 中动作的 ** 先验概率 ** $P(s,a)$。

\subsubsection{价值网络 $V_\phi(s)$}

价值网络输出一个标量，预测从状态 $s$ 出发的期望累计奖励：
\[
  V_\phi(s)=\operatorname{Tanh}(f_{\mathrm{value}}(h_s)),
\]
其中 $f_{\mathrm{value}}$ 为另一组全连接层，输出经 $\tanh$ 激活以限制在 $[-1,1]$。该值用
于 MCTS 中的节点评估和提前终止判断。

\newpage

\subsection{蒙特卡洛树搜索（MCTS）}

MCTS 在每次决策时进行多步模拟，以平衡“探索”与“利用”。其核
心流程包含四个阶段：

\begin{algorithm}[H]
\caption{MCTS 搜索过程}
\KwInput{当前状态 $s_0$，策略网络 $\pi_\theta$，价值网络 $V_\phi$，模拟次数 $N_{\mathrm{sim}}$}
\KwOutput{改进后的动作概率分布 $\pi$（基于访问次数）}
初始化搜索树，根节点为 $s_0$\;
\For{$i=1$ \KwTo $N_{\mathrm{sim}}$}{
  $node \leftarrow$ 根节点\;
  \While{$node$ 非叶节点}{
    \[
      a \leftarrow \arg\max_{a'}
      \left(
      Q(node,a')+
      c \cdot P(node,a') \cdot
      \frac{\sqrt{\sum_b N(node,b)}}{1+N(node,a')}
      \right)
    \]
    \tcp{选择阶段}
    $node \leftarrow$ 子节点$(node,a)$\;
  }
  \uIf{$node$ 为终止状态}{
    $z \leftarrow R_{\mathrm{final}}$\;
  }
  \Else{
    \For{每个动作 $a'$ 在动作空间中}{
      创建新子节点 $child$，其先验 $P(child)=\pi_\theta(a'|s_{node})$\;
    }
    $z \leftarrow$ 模拟$(node)$\tcp*[r]{使用默认策略快速走完}
  }
  \tcp{反向传播}
  \While{$node \ne$ 根节点}{
    更新 $N(node,a) += 1$，$Q(node,a) += z - V_\phi(s_{node})$\;
    $node \leftarrow$ 父节点$(node)$\;
  }
}
\Return{$\pi(a|s_0)=N(s_0,a)/\sum_b N(s_0,b)$}\tcp*[r]{访问次数归一化}
\end{algorithm}

其中 $c$ 为探索常数，$N$ 为访问次数，$Q$ 为平均价值，模拟阶段可采用
随机走子或轻量级启发式策略。

\newpage

\subsection{训练循环}

整个系统通过“搜索-存储-学习”的闭环不断进化。训练算法如算
法~?? 所示。

\begin{algorithm}[H]
\caption{训练循环}
\KwInput{布尔函数数据集 $\mathcal{D}$，策略网络 $\pi_\theta$，价值网络 $V_\phi$，训练轮数 $E$}
\KwOutput{训练好的网络参数 $\theta,\phi$}
初始化经验池 $\mathcal{B}\leftarrow\emptyset$\;
\For{$epoch=1$ \KwTo $E$}{
  \ForEach{函数 $f \in \mathcal{D}$}{
    使用 MCTS（算法 1）搜索得到改进策略 $\pi$ 和最终奖励 $z$\;
    记录每一步的 $(s_t,\pi_t,z_t)$ 存入 $\mathcal{B}$\;
  }
  \For{$step=1$ \KwTo $K$}{
    从 $\mathcal{B}$ 中采样一批数据 $(s,\pi,z)$\;
    计算损失：
    \[
      \mathcal{L}
      =
      \frac{1}{|\mathrm{batch}|}
      \sum_i
      \left[
        (z_i-V_\phi(s_i))^2
        - \pi_i^\top \log \pi_\theta(s_i)
      \right]
      + \lambda\|\theta\|_2^2,
    \]
    其中 $\lambda$ 为正则系数。使用 Adam 优化器更新 $\theta,\phi$\;
  }
}
\Return{$\theta,\phi$}
\end{algorithm}

\subsection{动作空间与门库}

动作空间 $\mathcal{A}$ 可根据研究目标灵活配置。本文考虑以下常用门库：

\begin{itemize}
  \item \textbf{CNT 门库：} $\{NOT,CNOT,Toffoli\}$，适合可逆逻辑综合；

  \item \textbf{通用门库：} $\{H,X,Y,Z,CNOT,T,Toffoli,SWAP\}$，表达能力最强；

  \item \textbf{Clifford+T 门库：} $\{H,S,CNOT,T\}$，贴合容错计算硬件约束。
\end{itemize}

\newpage

每个动作同时包含门类型和作用量子比特（如 $CNOT_{0,1}$），动作空间大小为
$|\text{门库}| \times n(n-1)/2$（对于双比特门）或 $|\text{门库}| \times n$（对于单比特门）。

\subsection{复杂性分析}

MCTS 每次模拟的复杂度为 $O(L\cdot|\mathcal{A}|)$，其中 $L$ 为模拟深度，$|\mathcal{A}|$ 为动
作空间大小；训练过程的计算成本主要取决于模拟次数和数据规模。由于策
略网络和价值网络提供了先验和评估，实际搜索效率远优于穷举。

\section{小结}

本章详细阐述了基于 RL 与 MCTS 的量子布尔函数综合方法，包括
MDP 建模、双网络架构、MCTS 搜索算法及训练流程。该方法通过智能体
的自我博弈，能够在指数级搜索空间中高效地生成高质量量子电路，为后续
实验验证奠定了理论基础。

\end{document}
